
Each of the once-through transition scenarios are compared on multiple 
criteria: the energy supplied, the number of advanced reactors deployed, 
the mass of enriched 
uranium required, the amount of \gls{SWU} capacity required to enrich uranium,
and the mass of waste produced. A subset of these results are presented
in \cite{bachmann_enrichment_2021}, except the \glspl{LWR} were assumed to 
operate for 60 years if not closed before December 2020. The results of Scenario 
1 are presented first, followed by the results of the no growth scenarios 
(Scenarios 2-7), then the results of the 1\% growth scenarios (Scenarios 8-13)
for each metric. 

\section{Scenario 1}
Scenario 1 models only the \glspl{LWR} deploying the United States with no 
prescribed energy demand. This scenario provides insight on historical 
requirements of the nuclear indsutry as well as future demands if no 
new reactors are deployed. Figure \ref{fig:energy_reactor1} shows the number of 
reactors deployed in Scenario 1 as a function of time comapred to the 
energy produced by the \glspl{LWR} each year. The engery produced by the 
\glspl{LWR} follows with the number of reactors that are deployed. A maximum of 109 
\glspl{LWR} are deployed at one time in this scenario, producing  and 92 \glspl{LWR}
are deployed when the transition begins in 2025. The first \gls{LWR} is 
commissioned in August 1967All of the \glspl{LWR} are
decommissioned by October of 2055. In 2025, the \glspl{LWR} produce 
89.45 GWe-y of energy, which is used to define the energy demand for 
the other scenarios. 

\begin{figure}
    \centering
    \includegraphics[scale=0.8]{s1_energy_reactors.pdf}
    \caption{Energy supplied by \glspl{LWR} in Scenario 1 each year of
    the simulation compared to the number of reactors deployed.}
    \label{fig:energy_reactor1}
\end{figure}


The mass of enriched uranium sent to the \glspl{LWR} in this scenario (Figure 
\ref{fig:fuel1}) follows the general pattern of the reactor deployment, but with 
more variation between individual time steps. The increased variation between 
time steps is becuase of the staggering of outages for the reactors, or when 
they receive fuel. There are also times with large increases in the mass of fuel 
sent to the reactors, such as in 2016, because these times include the deployment 
of a new reactor with a full core of fuel instead of the amount that is 
sent for refueling. 

The maximum amount of enriched uranium sent to the \glspl{LWR} at any one 
time in this scenario is 513.72 MTU. The average mass of enriched uranium sent to 
the reactors for the entire scenario is 135.74 MTU/month. The average mass of 
enriched uranium sent to the \glspl{LWR} between when they are first deployed 
and 2025 is 164.92 MTU/month. This average is larger than what is reported in 
\cite{bachmann_enrichment_2021}, and is because the average reported by 
Bachmann et. al accounts for the time before the \glspl{LWR} are actually 
deployed while the average reported here does not. The average mass of 
uranium sent to \glspl{LWR} between 2025-2055, from the start of the transition 
to when they are all decommissioned is 81.11 MTU/month. Comparing the averages 
shows that the vast majority of enriched uranium is sent to the \glspl{LWR} 
before 2025, which matches with the decline in the number of \glspl{LWR} 
deployed after 2020. After 2025, a total of 29,848.6 MTU is sent to the \glspl{LWR},
which is less than the cumulative mass reported in \cite{bachmann_enrichment_2021}
because of the different assumption of when the \glspl{LWR} will 
be decommissioned. 

\begin{figure}
    \centering
    \includegraphics[scale=0.8]{s1_uox.pdf}
    \caption{Mass of uranium supplied to the LWRs in Scenario 1 at each time step.}
    \label{fig:fuel1}
\end{figure}

The next metric of interest for this work is the mass of natural uranium 
required as feed material to produce the enriched uranium for the 
reactors. The mass of natural uranium required to produce the fuel sent to the 
reactors at 
is about an order of magnitude larger than the mass of the enriched uranium 
(Figure \ref{fig:feed1}). A maximum of 4,121.8 MTU of natural uranium is needed 
at one time in Scenario 1. While the \glspl{LWR} are deployed, an average of 
1,089.1 MTU/month of natural uranium is needed. Before 2025, an average of 
1,323.2 MTU/month of natural uranium is needed and after 2025 the average is 
650.8 MTU/month. 

\begin{figure}
    \centering
    \includegraphics[scale=0.8]{s1_feed.pdf}
    \caption{Mass of natural uranium required to supply the fuel sent to LWRs at each 
    time step in Scenario 1.}
    \label{fig:feed1}
\end{figure}

The next metric of interest for this work is the \gls{SWU} capacity required 
to produce the enriched uranium for the \glspl{LWR}. This result is 
shown in Figure \ref{fig:swu1}. A maximum of 3.71$\times 10^6$ kg-SWU is required to 
enrich the uranium sent to the reactors at one time step. An averge of 
0.981$\times 10^6$ kg-SWU/month is needed to enrich all of the uranium sent to 
the \glspl{LWR} during the time the \glspl{LWR} are deployed. An averge of 
1.19$\times10^6$ kg-SWU/month is required to produce the enriched uranium 
for \glspl{LWR} between their initial deployment and 2025, and 
0.586$\times 10^6$ kg-SWU/month is needed to produce the 
enriched uranium sent to \glspl{LWR} between 2025 and 2055. The values 
reported here are each different that what is reported in \cite{bachmann_enrichment_2021}
for a few different reasons. The first is that these calculations only include
the time steps in which the \glspl{LWR} are deployed while Bachmann et. al 
does not, similar to what leads to the differences in the reported mass of 
enriched uranium. The other factor contributing to the different values 
is the use of different enrichment levels for the \gls{LWR} fuel. 
Bachmann et. al uses an enrichment of 4.5\% while this work uses 4.3\%. 


\begin{figure}
    \centering
    \includegraphics[scale=0.8]{s1_swu.pdf}
    \caption{SWU capacity required to enrich the uranium sent to LWRs at each time step in Scenario 1.}
    \label{fig:swu1}
\end{figure}

Finally, the waste discharged from the reactors as a function of time is shown 
in Figure \ref{fig:waste1}. An maximum of 393.6 MT of spent fuel is discharged 
from the \glspl{LWR} at one time step. An average of 130.1 MT/month of spent fuel is 
discharged from the reactors during the time steps with \glspl{LWR} deployed.
An average of 149.5 MTU/month is discharged between the first \gls{LWR} deployment 
and 2025, an average of 93.9 MT/month of spent fuel is discharged from reactors 
between 2025-2055. 

\begin{figure}
    \centering
    \includegraphics[scale=0.8]{s1_waste.pdf}
    \caption{Mass of spent fuel dishcarged from the reactors in Scenario 1 as a function of time.}
    \label{fig:waste1}
\end{figure}

\section{Energy Supplied}
\subsection{No growth scenarios}
\begin{figure}
    \centering
    \includegraphics{nogrowth_energy.pdf}
    \caption{Energy supplied by reactors in each scenario, compared to 
    the energy demand (the black line)}
    \label{fig:energy_nogrowth}
\end{figure}

\subsection{1\% growth scenarios}
\begin{figure}
    \centering
    \includegraphics{1percent_energy.pdf}
    \caption{Energy supplied by reactors in each scenario, compared to 
    the energy demand (the black line)}
    \label{fig:energy_1percent}
\end{figure}

\section{Reactor deployment}
Each of the scenarios are compared on the number of advanced reactors that must 
be built to meet the prescribed energy demand. 

The \Cycamore `ManagerInst' has preferential deployment of prototypes based on 
the order in which they are listed in the input file. The first prototype 
is deployed until deploying more would create a surplus of the required 
commodity (in this case power), then the next prototype is deployed until 
its deployment would create a surplus of the requried commodity, and so forth. 
Based on this, the prototypes were listed in decreasing order of power 
output (i.e. Xe-100, VOYGR, then \gls{MMR}, for the necessary prototypes 
in that scenario) so that the number of reactors deployed is optimized to 
be the fewest reactors required for that scenario. 
\subsection{No growth scenarios}

\begin{figure}
    \centering
    \includegraphics[scale=0.8]{nogrowth_reactors.pdf}
    \caption{Total number of advanced reactors as a function of time in Scenarios 2-7.}
    \label{fig:nogrowth_reactors}
\end{figure}

\begin{table}
    \centering 
    \caption{Maximum number of each type of advanced reactor and maximum total 
    number of reactors deployed in Scenarios 2-7}
    \label{tab:reactors_nogrowth}
    \begin{tabular}{c c c c c}
        \hline
        Scenario & \glspl{MMR} & Xe-100s & VOYGRs & Total\\\hline
        2 & 9182 & - & - & 9182\\
        3 & - & 1225 & - & 1225\\
        4 & 752 & 1124 & - & 1876\\
        5 & 127 & - & 1811 & 1938\\
        6 & - & 1099 & 191 & 1288\\
        7 & 782 & 1122 & 0 & 1902\\
        \hline
    \end{tabular}
\end{table}

\begin{table}
    \centering 
    \caption{Maximum number of each type of advanced reactor and maximum total 
    number of reactors added to the simulation at a single time step in Scenarios 2-7.}
    \label{tab:reactors_added_nogrowth}
    \begin{tabular}{c c c c c}
        \hline
        Scenario & \glspl{MMR} & Xe-100s & VOYGRs & Total\\\hline
        2 & 378 & - & - & 378\\
        3 & - & 50 & - & 50\\
        4 & 31 & 50 & - & 56\\
        5 & 4 & - & 75 & 78\\
        6 & - & 49 & 8 & 51\\
        7 & 30 & 50 & 0 & 56\\
        \hline
    \end{tabular}
\end{table}

\subsection{1\% growth scenarios}
\begin{figure}
    \centering
    \includegraphics[scale=0.8]{1percent_reactors.pdf}
    \caption{Total number of advanced reactors as a function of time in Scenarios 8-13.}
    \label{fig:1percent_reactors}
\end{figure}
\begin{table}
    \centering 
    \caption{Maximum number of each type of advanced reactor and maximum total 
    number of reactors deployed in Scenarios 8-13}
    \label{tab:reactors_1growth}
    \begin{tabular}{c c c c c}
        \hline
        Scenario & \glspl{MMR} & Xe-100s & VOYGRs & Total\\\hline
        8 & \\
        9 & \\
        10 & \\
        11 & \\
        12 & \\
        13 & \\
        \hline
    \end{tabular}
\end{table}

\section{Uranium resources}

\subsection{No growth scenarios}

\begin{figure}
    \centering
    \begin{subfigure}[b]{0.45\textwidth}
        \centering
        \includegraphics[width=\textwidth]{nogrowth_uranium.pdf}
        \caption{Mass of uranium sent to all reactors at each time step.}
        \label{fig:nogrowth_all_urnaium}
    \end{subfigure}
    \hfill
    \begin{subfigure}[b]{0.45\textwidth}
        \centering
        \includegraphics[width=\textwidth]{nogrowth_AR_uranium.pdf}
        \caption{Mass of enriched uranium sent to advanced reactors at each time step.}
        \label{fig:nogrowth_AR_uranium}
    \end{subfigure}
       \caption{Total mass enriched uranium required by reactors
       at each time step in Scenarios 2-7.}
       \label{fig:nogrowth_urnaium}
\end{figure}

\begin{figure}
    \centering
    \begin{subfigure}[b]{0.45\textwidth}
        \centering
        \includegraphics[width=\textwidth]{nogrowth_feed.pdf}
        \caption{Mass of feed uranium to produce enriched uranium sent to 
        all reactors at each time step.}
        \label{fig:nogrowth_all_feed}
    \end{subfigure}
    \hfill
    \begin{subfigure}[b]{0.45\textwidth}
        \centering
        \includegraphics[width=\textwidth]{nogrowth_AR_feed.pdf}
        \caption{Mass of feed uranium to produce enriched uranium sent to 
        advanced reactors at each time step.}
        \label{fig:nogrowth_AR_feed}
    \end{subfigure}
       \caption{Total mass of feed uranium required to produce enriched uranium
       at each time step in Scenarios 2-7.}
       \label{fig:nogrowth_feed}
\end{figure}

\subsection{1\% growth scenarios}

\begin{figure}
    \centering
    \begin{subfigure}[b]{0.45\textwidth}
        \centering
        \includegraphics[width=\textwidth]{1percent_uranium.pdf}
        \caption{Mass of uranium sent to all reactors at each time step.}
        \label{fig:1percent_all_urnaium}
    \end{subfigure}
    \hfill
    \begin{subfigure}[b]{0.45\textwidth}
        \centering
        \includegraphics[width=\textwidth]{1percent_AR_uranium.pdf}
        \caption{Mass of enriched uranium sent to advanced reactors at each time step.}
        \label{fig:1percent_AR_uranium}
    \end{subfigure}
       \caption{Total mass enriched uranium required by reactors
       at each time step in Scenarios 8-13.}
       \label{fig:1percent_urnaium}
\end{figure}

\begin{figure}
    \centering
    \begin{subfigure}[b]{0.45\textwidth}
        \centering
        \includegraphics[width=\textwidth]{1percent_feed.pdf}
        \caption{Mass of feed uranium to produce enriched uranium sent to 
        all reactors at each time step.}
        \label{fig:1percent_all_feed}
    \end{subfigure}
    \hfill
    \begin{subfigure}[b]{0.45\textwidth}
        \centering
        \includegraphics[width=\textwidth]{1percent_AR_feed.pdf}
        \caption{Mass of feed uranium to produce enriched uranium sent to 
        advanced reactors at each time step.}
        \label{fig:1percent_AR_feed}
    \end{subfigure}
       \caption{Total mass of feed uranium required to produce enriched uranium
       at each time step in Scenarios 8-13.}
       \label{fig:1percent_feed}
\end{figure}
\section{SWU capacity}
\subsection{No growth scenarios}

\begin{figure}
    \centering
    \begin{subfigure}[b]{0.45\textwidth}
        \centering
        \includegraphics[width=\textwidth]{nogrowth_SWU.pdf}
        \caption{\gls{SWU} capacity required to enrich the mass of 
        uranium sent to reactors at each time step.}
        \label{fig:nogrowth_all_SWU}
    \end{subfigure}
    \hfill
    \begin{subfigure}[b]{0.45\textwidth}
        \centering
        \includegraphics[width=\textwidth]{nogrowth_AR_SWU.pdf}
        \caption{\gls{SWU} capacity required to enrich the mass of 
        uranium sent to advanced reactors at each time step.}
        \label{fig:nogrowth_AR_SWU}
    \end{subfigure}
       \caption{Total mass of fuel and enriched uranium required by all reactors 
       at each time step in Scenarios 2-7.}
       \label{fig:nogrowth_swu}
\end{figure}

\subsection{1\% growth scenarios}
\begin{figure}
    \centering
    \begin{subfigure}[b]{0.45\textwidth}
        \centering
        \includegraphics[width=\textwidth]{1percent_SWU.pdf}
        \caption{\gls{SWU} capacity required to enrich the mass of 
        uranium sent to reactors at each time step.}
        \label{fig:1percent_all_SWU}
    \end{subfigure}
    \hfill
    \begin{subfigure}[b]{0.45\textwidth}
        \centering
        \includegraphics[width=\textwidth]{1percent_AR_SWU.pdf}
        \caption{\gls{SWU} capacity required to enrich the mass of 
        uranium sent to advanced reactors at each time step.}
        \label{fig:1percent_AR_SWU}
    \end{subfigure}
       \caption{Total mass of fuel and enriched uranium required by all reactors 
       at each time step in Scenarios 8-13.}
       \label{fig:1percent_swu}
\end{figure}

\section{Waste}
\subsection{No growth scenarios}

\begin{figure}
    \centering
    \begin{subfigure}[b]{0.45\textwidth}
        \centering
        \includegraphics[width=\textwidth]{nogrowth_waste.pdf}
        \caption{Mass of fuel discharged from all reactors as 
        a function of time.}
        \label{fig:nogrowth_all_waste}
    \end{subfigure}
    \hfill
    \begin{subfigure}[b]{0.45\textwidth}
        \centering
        \includegraphics[width=\textwidth]{nogrowth_AR_waste.pdf}
        \caption{Mass of fuel discharged from advanced reactors as 
        a function of time..}
        \label{fig:nogrowth_AR_waste}
    \end{subfigure}
       \caption{Mass of fuel discharged from reactors 
       as a function of time for Scenarios 2-7. }
       \label{fig:nogrowth_waste}
\end{figure}

\subsection{1\% growth scenarios}
\begin{figure}
    \centering
    \begin{subfigure}[b]{0.45\textwidth}
        \centering
        \includegraphics[width=\textwidth]{1percent_waste.pdf}
        \caption{Mass of fuel discharged from all reactors as 
        a function of time.}
        \label{fig:1percent_all_waste}
    \end{subfigure}
    \hfill
    \begin{subfigure}[b]{0.45\textwidth}
        \centering
        \includegraphics[width=\textwidth]{1percent_AR_waste.pdf}
        \caption{Mass of fuel discharged from advanced reactors as 
        a function of time..}
        \label{fig:1percent_AR_waste}
    \end{subfigure}
       \caption{Mass of fuel discharged from reactors 
       as a function of time for Scenarios 8-13. }
       \label{fig:1percent_waste}
\end{figure}